\chapter{Методика}
\label{ch:methods}

Узнать тип и концентрацию носителей заряда, проводимость материала и подвижность носителей в полупроводниках можно с помощью постоянной Холла.

Концентрацию можно найти с помощью \cite{Lab}:
\begin{equation}\label{equ:R_H}
    n = -\dfrac{1}{R_H e}
\end{equation}
где $n$ — концентрация носителей заряда, м$^{-3}$; $R_H$ — постоянная Холла, м$^{3}$/Кл;
$e = 1{,}60 \cdot 10^{-19}$ Кл — модуль заряда электрона.

По знаку постоянной Холла можно определить тип носителей заряда (знак «–» для n-типа проводимости, «+» - для p-типа).

Подвижность носителей можно вычислить с помощью \cite{Lab}:
\begin{equation}
    \mu = \sigma |R_H| = \frac{|R_H|}{\rho},
    \label{eq:mu}
\end{equation}
где $\mu$~— подвижность носителей, м$^{2}$/(В·с); $\sigma$~— удельная проводимость, (Ом·м)$^{-1}$; $\rho$~— удельное сопротивление, Ом·м.

Уделььную проводимость можно найти по формуле \cite{Lab}:
\begin{equation}
    \sigma = n e \mu \quad (n\text{-тип}), \qquad 
    \sigma = p e \mu \quad (p\text{-тип}),
    \label{eq:sigma}
\end{equation}
где $\sigma$~— удельная проводимость, (Ом·м)$^{-1}$; $n$~— концентрация электронов, м$^{-3}$; $p$~— концентрация дырок, м$^{-3}$; $e = 1,60 \cdot 10^{-19}$ Кл~— элементарный заряд; $\mu$~— подвижность носителей, м$^{2}$/(В·с).

Найти постоянную Холла можно при помощи измерения разности потенциалов $\Delta E$ между холловскими контактами при 
фиксированном приращении магнитного поля $\Delta B$ \cite{Lab}:

\begin{equation}
    \Delta E = U_H = R_H \cdot \frac{I \cdot \Delta B}{d},
\end{equation}
где $U_H$ — напряжение Холла, $I$ — ток через образец, $\Delta B$ — приращение магнитного поля, $d$ — толщина образца.

Если зависимости окажутся линейными, то коэффициенты наклона К при фиксированных токах \cite{Lab}:
\begin{equation}
    K = \frac{\Delta E}{\Delta B} = \left( \frac{R_H}{d} \right) \cdot I.
\end{equation}

Построив график K(I), если зависимость также окажется линейной, постоянную Холла можно найти при помощи \cite{Lab}:

\begin{equation}
    R_H = \frac{K d}{I}.
\end{equation}
где $R_H$ - постоянная Холла; $\frac{K}{I}$ - коэффициент наклона прямой; $d$ - толщина исследуемого образца.